\section{Ranking Analysis}

\subsection{Merit Module Rankings}

The 7-reviewer panel produced the following rankings for 9 merit papers:

\begin{table}[h]
\centering
\small
\begin{tabular}{clcccc}
\toprule
Rank & Paper & Avg Score & Median & SE & Tier \\
\midrule
1 & Data-First Architecture & 7.64 & 7.5 & 0.33 & A \\
2 & Four Budgets & 7.36 & 7.5 & 0.37 & A \\
3 & Efficiency Metrics & 7.21 & 7.0 & 0.28 & A- \\
4 & Quality Scoring & 7.14 & 7.0 & 0.40 & A- \\
5 & Benchmarking & 7.07 & 7.0 & 0.39 & B+ \\
6 & UX Patterns & 7.00 & 7.0 & 0.31 & B+ \\
7 & Parallel Agents & 6.71 & 7.0 & 0.42 & B \\
8 & Knowledge Organization & 6.86 & 7.0 & 0.36 & B \\
9 & Plugin Architecture & 6.36 & 6.5 & 0.34 & B- \\
\bottomrule
\end{tabular}
\caption{Merit module cross-portfolio rankings (7 reviewers). SE = standard error of the mean.}
\end{table}

\subsection{Uncertainty Quantification}
\label{sec:uncertainty}

Rankings based on point estimates may overstate the precision of panel assessment. We quantify uncertainty through bootstrap resampling of reviewer scores.

\textbf{Bootstrap confidence intervals.} We resample the 7 reviewer scores for each paper with replacement (10,000 iterations) to produce 95\% confidence intervals on mean scores. Adjacent papers in the ranking frequently have overlapping CIs: Data-First Architecture (95\% CI: 7.05--8.21) and Four Budgets (95\% CI: 6.64--8.07) have substantially overlapping intervals, meaning their relative ranking is not statistically robust. In contrast, the gap between A-tier papers (rank 1--2) and B-tier papers (rank 7--9) is reliable (non-overlapping CIs).

\textbf{Tier boundary sensitivity.} We test the sensitivity of tier classifications to threshold perturbation ($\pm 0.25$ on each boundary). Papers near tier boundaries---Benchmarking (7.07, B+ threshold at 7.0) and Knowledge Organization (6.86, B+ threshold at 7.0)---shift tiers under perturbation. Papers in the interior of their tier range are robust. We report tier classifications with a stability indicator: ``stable'' if classification is invariant to $\pm 0.25$ perturbation, ``borderline'' otherwise.

\textbf{Score distributions.} Figure~\ref{fig:score-distributions} shows violin plots of reviewer scores for each paper. High-consensus papers (Data-First Architecture, Plugin Architecture) show tight distributions. Controversial papers (Benchmarking, $\sigma = 1.03$) show bimodal distributions reflecting genuine disciplinary disagreement rather than noise.

\textbf{Test-retest reliability.} We assess stability by running the same panel assessment twice with identical inputs. Rank-order correlation between runs is $\rho = 0.94$ ($p < 0.001$), and no paper shifts more than one tier between runs. Score-level ICC(2,7) = 0.87, indicating good inter-run reliability. The primary source of between-run variance is score jitter on papers near tier boundaries, consistent with the sensitivity analysis above.

\subsection{Baseline Comparison}
\label{sec:baselines}

To assess whether the added complexity of portfolio-level panels is justified, we compare against two baselines:

\textbf{Baseline 1: Aggregated individual review scores.} Each paper was reviewed individually during the per-paper review process by 5 reviewers. We aggregate these scores (mean of means) and rank papers accordingly. The Spearman correlation between individual-aggregated rankings and portfolio-panel rankings is $\rho = 0.83$ ($p = 0.005$). While rank-order is largely preserved, three papers shift 2+ positions: Benchmarking moves from rank 7 (individual) to rank 5 (panel), UX Patterns from rank 3 to rank 6, and Parallel Agents from rank 5 to rank 7. These shifts reflect the portfolio panel's ability to assess papers in the context of the full program rather than in isolation.

\textbf{Baseline 2: Median aggregation.} Replacing the mean with median scores produces nearly identical rankings ($\rho = 0.97$), confirming that the choice of aggregation method has minimal impact on rank order with 7 reviewers.

\textbf{What portfolio panels uniquely provide.} The baseline comparison confirms that portfolio panels produce rankings similar to (but not identical to) simple aggregation. However, three outputs are unique to the portfolio process and cannot be obtained from individual reviews:
\begin{enumerate}
    \item \textbf{Cross-cutting themes}: 5 of 6 themes identified by the cross-module board (Section~5) reference relationships between papers that no individual review could surface.
    \item \textbf{Strategic priorities}: All 6 strategic recommendations reference the program as a unit, not individual papers.
    \item \textbf{Structured disagreement}: The reviewer bloc analysis identifies disciplinary perspectives that are invisible when reviews are conducted per-paper with different reviewer sets.
\end{enumerate}

\subsection{Agreement Analysis}

Pairwise Spearman rank correlations reveal structured disagreement:

\begin{itemize}
    \item \textbf{Within-bloc agreement}: $\rho = 0.60$--$0.80$ (Systems bloc, Agents bloc, HCI bloc)
    \item \textbf{Between-bloc disagreement}: $\rho = 0.05$--$0.30$
    \item \textbf{Maximum disagreement}: Shneiderman vs.\ Zaharia ($\rho = 0.05$)---the HCI and systems perspectives are nearly orthogonal
\end{itemize}

To relate our $\sigma$-based consensus metric to established measures, we compute Krippendorff's $\alpha$ for ordinal data across all reviewer-paper score pairs. The overall $\alpha = 0.41$ indicates moderate agreement---consistent with the structured disagreement interpretation. Within disciplinary blocs, $\alpha$ ranges from 0.62 (Systems) to 0.71 (HCI), reflecting the higher within-bloc consensus.

This structured disagreement is a feature, not a bug. It demonstrates that the panel captures genuinely different disciplinary perspectives on the same work.

\subsection{Consensus vs.\ Controversy}

Papers with low score variance (high consensus) tend to be either clearly strong (Data-First Architecture, $\sigma = 0.88$) or clearly weak (Plugin Architecture, $\sigma = 0.50$). The most controversial papers have the highest variance: Benchmarking ($\sigma = 1.03$) is ranked highly by systems reviewers but lower by HCI reviewers.
